\pdfoutput=1
\documentclass[11pt]{article}
\usepackage[T1]{fontenc}
\usepackage{lmodern}
\usepackage{amsmath,amssymb,amsthm,mathtools}
\usepackage{booktabs}
% Generated by scripts/make_claims_tex.py from claims.json -- do not edit by hand.
\newcommand{\clDate}{November 21, 2026}
\newcommand{\clpdfEpoch}{1{,}795{,}219{,}200}
\newcommand{\clpages}{22}
\newcommand{\clcadicalVersion}{3.1.0}
\newcommand{\clpysatVersion}{1.9.dev24}
\newcommand{\clpairsTotal}{10{,}888}
\newcommand{\clpairsSix}{10{,}887}
\newcommand{\clsurvSha}{58068a1dd390e350}
\newcommand{\clvNodes}{31{,}191{,}558{,}503}
\newcommand{\clvCpu}{15{,}019}
\newcommand{\clsatSolve}{1137}
\newcommand{\clsatCheck}{1498}
\newcommand{\clsatMonoSolve}{221}
\newcommand{\clsatMonoCheck}{408}
\newcommand{\clsatOtherMax}{3}
\newcommand{\clhStageOne}{7.0}
\newcommand{\clhStageTwo}{2.6}
\newcommand{\clhStageOneV}{4.2}
\newcommand{\clhSat}{0.7}
\newcommand{\clhCert}{2.0}
\newcommand{\clcertProofGB}{25}
\newcommand{\clcubesSix}{1072}
\newcommand{\clcubesSmall}{76}
\newcommand{\clsampleParts}{31}
\newcommand{\clbruteMaxN}{12}
\newcommand{\clrefCompareMaxN}{18}
\newcommand{\clrefPipeMaxN}{12}
\newcommand{\clrefDirectMaxN}{10}
\newcommand{\clpairsTotalXXVIII}{118{,}765}
\newcommand{\clpairsSevenXXVIII}{118{,}662}
\newcommand{\clpairsSixXXVIII}{101}
\newcommand{\clsurvShaXXVIII}{7bad471036c300f2}
\newcommand{\clpartsXXVIII}{16{,}800}
\newcommand{\clcloudPartsXXVIII}{7457}
\newcommand{\clcloudCallsXXVIII}{933}
\newcommand{\clsatSolveXXVIII}{21{,}985}
\newcommand{\clsatCheckXXVIII}{24{,}017}
\newcommand{\clsatMonoSolveXXVIII}{1458}
\newcommand{\clsatMonoCheckXXVIII}{4452}
\newcommand{\clsatOtherMaxXXVIII}{32}
\newcommand{\clhStageOneXXVIII}{83.1}
\newcommand{\clhSatXXVIII}{12.8}
\newcommand{\clhCertXXVIII}{26.5}
\newcommand{\clhStageOneVXXVIII}{4.2}
\newcommand{\clhPlaceXXVIII}{0.4}
\newcommand{\clcertProofGBXXVIII}{182}
\newcommand{\clcertMaxProofGBXXVIII}{1.26}
\newcommand{\clclausesSevenXXVIII}{477{,}629}
\newcommand{\clvPartsSixXXVIII}{1549}
\newcommand{\clvPairsSixXXVIII}{90}
\newcommand{\clvNodesXXVIII}{38{,}258{,}356{,}242}
\newcommand{\clvCpuXXVIII}{15{,}066}
\newcommand{\clplacePairsXXVIII}{102}
\newcommand{\clplaceCpuXXVIII}{1501}
\newcommand{\clsamplePartsXXVIII}{20}
\newcommand{\clframesXXXV}{5050}
\newcommand{\clcasesXXXV}{128{,}941}
\newcommand{\clrefPlainXXXV}{128{,}356}
\newcommand{\clrefUnsplitXXXV}{508}
\newcommand{\clrefSplitXXXV}{77}
\newcommand{\clsecondLayerXXXV}{1}
\newcommand{\clcubesXXXV}{1755}
\newcommand{\cltimeoutRecXXXV}{8}
\newcommand{\clphaseOneCasesXXXV}{115{,}601}
\newcommand{\clphaseOneSolveXXXV}{16.4}
\newcommand{\clphaseOneCheckXXXV}{9.3}
\newcommand{\clphaseTwoCasesXXXV}{12{,}755}
\newcommand{\clphaseTwoSolveXXXV}{15.2}
\newcommand{\clphaseTwoCheckXXXV}{7.6}
\newcommand{\clphaseThreeCasesXXXV}{585}
\newcommand{\clphaseThreeSolveXXXV}{21.5}
\newcommand{\clphaseThreeCheckXXXV}{10.7}
\newcommand{\clhSolveXXXV}{53.1}
\newcommand{\clhCheckXXXV}{27.6}
\newcommand{\clproofGBXXXV}{232}
\newcommand{\clmaxProofGBXXXV}{1.25}
\newcommand{\clmaxPlainXXXV}{368}
\newcommand{\clmaxCaseXXXV}{3804}
\newcommand{\clrehashSample}{2000}
\newcommand{\clrehashSampleEven}{300}
\newcommand{\clrehashAllXXXV}{130{,}111}
\newcommand{\clresolveXXXV}{8}
\newcommand{\clresolveEven}{3}
\newcommand{\clleanTheorems}{116}
\newcommand{\clleanPieces}{343}
\newcommand{\clxKmMaxN}{35}
\newcommand{\clxKmCasesXXXV}{330}
\newcommand{\clxAllThirteenCases}{5103}
\newcommand{\clxAllThirteenSat}{26}
\newcommand{\clxAllFifteen}{7916}
\newcommand{\clxHardCases}{585}
\newcommand{\clxSampleCases}{2200}
\newcommand{\clxHardSplit}{14}
\newcommand{\clxHardLeaves}{795}
\newcommand{\clxHardRehash}{809}
\newcommand{\clxHardCpuH}{76.5}
\newcommand{\clxSampleCpu}{50{,}142}
\newcommand{\clrowZa}{$(0,13)$ & 792\,319 & 1 & 0 & 9\,367\,246\,608 & 7132 & 0 & 76 & 3\\}
\newcommand{\clrowZb}{$(1,12)$ & 3\,426\,552 & 0 & 0 & -- & -- & -- & 76 & 3\\}
\newcommand{\clrowZc}{$(2,11)$ & 15\,749\,727 & 0 & 5 & -- & -- & -- & 76 & 5\\}
\newcommand{\clrowZd}{$(3,10)$ & 75\,767\,491 & 0 & 27 & -- & -- & -- & 76 & 15\\}
\newcommand{\clrowZe}{$(4,9)$ & 546\,916\,450 & 0 & 256 & -- & -- & -- & 76 & 71\\}
\newcommand{\clrowZf}{$(5,8)$ & 4\,296\,637\,439 & 0 & 2416 & -- & -- & -- & 76 & 426\\}
\newcommand{\clrowZg}{$(6,7)$ & 29\,237\,608\,861 & 10\,887 & 22\,318 & 3\,408\,018\,489 & 2147 & 0 & 1072 & 6719\\}
\newcommand{\clrowYa}{$(0,14)$ & 1\,344\,512 & 1 & 0 & 76 & 3 & 5910\\}
\newcommand{\clrowYb}{$(1,13)$ & 4\,224\,030 & 1 & 0 & 76 & 3 & 0\\}
\newcommand{\clrowYc}{$(2,12)$ & 21\,160\,768 & 0 & 5 & 76 & 6 & --\\}
\newcommand{\clrowYd}{$(3,11)$ & 125\,674\,586 & 0 & 33 & 76 & 19 & --\\}
\newcommand{\clrowYe}{$(4,10)$ & 998\,940\,852 & 0 & 313 & 76 & 86 & --\\}
\newcommand{\clrowYf}{$(5,9)$ & 8\,199\,312\,979 & 0 & 3536 & 76 & 483 & --\\}
\newcommand{\clrowYg}{$(6,8)$ & 64\,418\,002\,306 & 101 & 46\,991 & 1072 & 5655 & 21\\}
\newcommand{\clrowYh}{$(7,7)$ & 251\,094\,656\,771 & 118\,662 & 248\,194 & 1072 & 88\,999 & 40\,071\\}
\newcommand{\clvalFour}{3}
\newcommand{\clvalSix}{1}
\newcommand{\clvalEight}{0}
\newcommand{\clvalTen}{1}
\newcommand{\clvalTwelve}{1}
\newcommand{\clvalFourteen}{0}
\newcommand{\clvalSixteen}{0}
\newcommand{\clvalEighteen}{2}
\newcommand{\clpairsrecTwenty}{23{,}888}
\newcommand{\clpairsrecTwentyTwo}{5087}
\newcommand{\clpairsrecTwentyFour}{56{,}957}
\newcommand{\clpairsgpEight}{871}
\newcommand{\clpairsgpFourteen}{571{,}200}
\newcommand{\clpairsgpSixteen}{555{,}156}
\newcommand{\clgpEight}{638}
\newcommand{\clgpFourteen}{588}
\newcommand{\clgpSixteen}{371}
\newcommand{\clrowFa}{26 & 13 & 1596 & 5663 & 413 & 549 & 0.59 & 17.8\\}
\newcommand{\clrowFb}{28 & 14 & 2080 & 7300 & 652 & 787 & 0.93 & 17.8\\}
\newcommand{\clrowFc}{32 & 16 & 3321 & 11\,466 & 1597 & 1559 & 2.12 & 29.6\\}
\newcommand{\clrowFd}{34 & 17 & 4095 & 14\,049 & 2548 & 2193 & 3.06 & 81.1\\}
\newcommand{\clrowFe}{36 & 18 & 5050 & 17\,205 & 3604 & 2927 & 4.32 & 57.5\\}
\newcommand{\clrowFg}{40 & 20 & 7381 & 24\,871 & 7892 & 5395 & 8.15 & 84.8\\}
\newcommand{\clrowFh}{42 & 21 & 8778 & 29\,447 & 11\,299 & 7154 & 10.85 & 132.7\\}
\newcommand{\clrowFi}{44 & 22 & 10\,440 & 34\,854 & 13\,243 & 8772 & 14.46 & 119.1\\}
\newcommand{\clrowFj}{46 & 23 & 12\,246 & 40\,728 & 10\,009 & 7724 & 18.55 & 135.3\\}
\newcommand{\clrowFf}{35 & 17 & 330 & 330 & 13 & 36 & 0.01 & 0.1\\}

\usepackage[margin=1in]{geometry}
\usepackage[colorlinks=true,urlcolor=blue,citecolor=blue,linkcolor=blue]{hyperref}

\newtheorem{theorem}{Theorem}[section]
\newtheorem{proposition}[theorem]{Proposition}
\newtheorem{lemma}[theorem]{Lemma}
\newtheorem{corollary}[theorem]{Corollary}
\theoremstyle{definition}
\newtheorem{definition}[theorem]{Definition}
\theoremstyle{remark}
\newtheorem{remark}[theorem]{Remark}

\emergencystretch=3em
\newcommand{\repo}{\url{https://github.com/jasondavies/queen-domination}}

\title{The domination numbers of the $26\times26$ and $28\times28$ queen's graphs}
\author{Jason Davies\thanks{\href{mailto:jason@jasondavies.com}{\nolinkurl{jason@jasondavies.com}}}}
\date{November 20, 2026}

\begin{document}

\maketitle
\begin{abstract}
We show that $14$ queens are needed to occupy or attack every square of a $26\times26$ chessboard and
$15$ queens for a $28\times28$ board, so $\gamma(Q_{26})=14$ and $\gamma(Q_{28})=15$. These were the two
smallest open cases of the queen domination problem (OEIS A075458); the previous bounds were
$13\le\gamma(Q_{26})\le14$ and $14\le\gamma(Q_{28})\le15$. Together with known results for $n=27$ and
$n=29,30,31$, the values of $\gamma(Q_n)$ are now known for all $n\le 31$. The proofs are
computer-assisted. If $n/2$ queens dominate the $n\times n$ board, $z_q$ of them on squares of colour
$q$, choose $n/2$ empty rows and $n/2$ empty columns. Every square of colour $q$ in both chosen sets is
attacked along a diagonal, so these squares are covered by at most $z_q$ difference diagonals and at
most $z_q$ sum diagonals of colour $q$. This is a constrained vertex cover problem in a small bipartite
graph, and a matching bound prunes it well enough that the possible pairs (empty rows, empty columns)
can be enumerated: up to symmetry there are $\clpairsTotal$ of them for $n=26$ and $\clpairsTotalXXVIII$
for $n=28$. A second search shows that no $n/2$ queens avoiding the chosen rows and columns dominate the
board, for any of these pairs. Both steps are certified by DRAT proofs checked with \texttt{drat-trim}:
the enumeration through a cube-and-conquer refutation of the relaxation with every listed pair
excluded, the second step by one refutation per pair. The lower bounds therefore depend on two short CNF
encoders and the proof checker, and not on the search programs. The same pipeline reproduces all known
values and isomorphism-class counts of minimum dominating sets of $Q_n$ for even $n\le 18$, and the
known values for $n=20,22,24$.
\end{abstract}

\noindent\textbf{AI transparency statement.} The result, the computer-assisted proof, the code and
this manuscript were produced by artificial intelligence. The human author initiated the broader
research and directed the AI systems, but did not write the text and has not checked the arguments.
Details are given at the end of the paper.

\section{Introduction}

The queen's graph $Q_n$ has the $n^2$ squares of the $n\times n$ board as vertices, two squares being
adjacent when they share a row, a column or a diagonal. Its domination number $\gamma(Q_n)$ is the least
number of queens that occupy or attack every square. The problem goes back to de~Jaenisch (1862) and
is Problem~C18 in Guy's collection; see the surveys \cite{Cockayne,Watkins} and the introduction
of~\cite{Weakley2022}. The general lower bound
\[
  \gamma(Q_n)\ \ge\ \tfrac{n-1}{2}
\]
was proved independently by Spencer (see~\cite{Cockayne}) and by Raghavan and Venkatesan~\cite{RV}, and
equality holds only for $n\in\{3,11\}$~\cite{Weakley2002,FW}. In particular $\gamma(Q_n)\ge n/2$ for even
$n$. For every $n\le120$ the value is known to be $\lceil n/2\rceil$ or $\lceil n/2\rceil+1$
\cite{OW,Weakley2022}. Exact values were known for $n\le 25$, for $n=27$, and for $n=29,30,31$
(Corollary~\ref{cor:oeis} below), as well as for some larger $n$. For $n\le25$ they come from constructions~\cite{OW,Kearse},
from exhaustive searches by Kearse and Gibbons ($n\le19$)~\cite{Kearse} and by Bird ($n=20,22,24$, with a
general-purpose dominating-set solver on 64 processors)~\cite{Bird}, and, for $n\le19$, from
proof-producing SAT solving~\cite{RB}. Bird identified $n=26$ as the smallest open case. A $14$-queen
dominating set of $Q_{26}$ was given by \"Osterg\aa rd and Weakley~\cite{OW}, who in fact found an
independent one, and another appears in Kamenetsky's list of best known placements attached to OEIS
A075458~\cite{OEIS}. So $13\le\gamma(Q_{26})\le 14$. The independent domination number of $Q_{26}$ is
already known to be $14$~\cite{LMR}. For $n=28$, \"Osterg\aa rd and Weakley showed
$i(Q_{28})\in\{14,15\}$~\cite{OW}, Langlois-R\'emillard, M\"ussig and Rold\'an computed $i(Q_{28})=15$~\cite{LMR},
and so $14\le\gamma(Q_{28})\le15$; once $n=26$ is settled, $n=28$ is the smallest open case.

For even $n=2k$ the question whether $\gamma(Q_n)=k$ is a tight one. Weakley~\cite{Weakley2022} calls a
dominating set \emph{monochromatic} if all its squares have the same colour. He shows that a
monochromatic $k$-set of $Q_{2k}$ forces $n\in\{4,12\}$ when $4\mid n$. When $n\equiv2\pmod 4$ it
forces a Diophantine condition, which has no solution for $n=26$. He also shows that the colour classes
of a bichromatic $k$-set have nearly equal sizes (for $k=14$: sizes $\{7,7\}$, or $\{6,8\}$ under further
conditions), and he conjectures that bichromatic $k$-sets do not exist for $n>4$. So the
open questions were whether $Q_{26}$ has a bichromatic $13$-dominating set and whether $Q_{28}$, where
monochromatic $14$-sets are excluded because $4\mid28$, has a bichromatic $14$-dominating set. Our
computation does not use these results: it treats all colour splits alike.

\begin{theorem}\label{thm:main}
$\gamma(Q_{26})=14$ and $\gamma(Q_{28})=15$. Equivalently, no $13$ queens dominate the $26\times26$
board, and no $14$ queens dominate the $28\times28$ board.
\end{theorem}

\begin{corollary}\label{cor:oeis}
The terms $a(n)=\gamma(Q_n)$ of OEIS A075458 are known for all $n\le31$:
\[
  a(26)=a(27)=14,\qquad a(28)=a(29)=a(30)=15,\qquad a(31)=16 .
\]
The smallest open case is now $n=32$, where $\gamma(Q_{32})\in\{16,17\}$.
\end{corollary}

\begin{proof}
By Corollary~2 of \"Osterg\aa rd and Weakley~\cite{OW}, $\gamma(Q_n)\ge n/2$ for $n<143$,
$n\ne3,11$; this gives the lower bounds $14$, $15$, $15$, $16$ for $n=27,29,30,31$. The upper bounds come
from explicit sets, whose domination our verification script checks directly.
For $n=27$ we use the $14$-set of \"Osterg\aa rd and Weakley \cite[\S5]{OW}: in centred coordinates
$-13,\dots,13$ it consists of the $13$ squares $(x,y)$ with $x=-12,-10,\dots,12$ and
$y=6,-12,-6,8,0,-8,10,-2,4,12,-10,-4,2$ respectively, together with $(-9,-9)$.
For $n=30$ we use Weakley's monochromatic $15$-set \cite[proof of Proposition~21]{Weakley2022}: in
coordinates $-14,\dots,15$, the squares $(x,y)$ with $x=-13,-11,\dots,15$ and
$y=1,-9,9,15,-11,-5,3,-3,13,11,7,-13,-7,-1,5$ respectively. None of its squares lies in the row or the
column with coordinate $-14$; deleting them leaves a $29\times29$ board on which the same $15$ squares
still dominate, so $\gamma(Q_{29})\le15$ (this value is also among those of~\cite{OW}). Adding a
$31$st row and column and a queen at their intersection gives a $16$-set dominating $Q_{31}$.
Finally $\gamma(Q_{32})\ge16$ and $\gamma(Q_{32})\le17$~\cite{Weakley2022}.
\end{proof}

Theorem~\ref{thm:main} also settles Weakley's Conjecture~4~\cite{Weakley2022} for $n=26$ and $n=28$,
vacuously.

Section~\ref{sec:relax} reduces the problem to a finite search, with complete proofs.
Section~\ref{sec:comp} describes the search and the two computations, and Section~\ref{sec:cert} the
DRAT certificates and what the proofs rest on. Section~\ref{sec:valid} reports the validation runs and the independent checks. The
code, the records of the computation and a verification script are available online.\footnote{\repo}

\section{A line relaxation}\label{sec:relax}

Index rows and columns by $[0,n)=\{0,\dots,n-1\}$ and write $(r,x)$ for the square in row $r$ and
column $x$. Its \emph{colour} is $r+x\bmod 2$. Its \emph{difference line} is $r-x$ and its \emph{sum
line} is $r+x$. Since $r-x\equiv r+x\pmod 2$, all squares of a line have the same colour, which we call
the colour of the line. A queen on $(r,x)$ lies on exactly one difference line and one sum line, both of
its own colour, and attacks diagonally only squares of its own colour.

Let $D$ be a dominating set of $Q_n$ with $|D|=g$, and let $z_0,z_1$ be the numbers of queens of
colour $0$ and $1$. Put $s=n-g$. Since $D$ meets at most $g$ rows and at most $g$ columns, at least $s$
rows and at least $s$ columns contain no queen.

\begin{definition}
For $A,B\subseteq[0,n)$ and $q\in\{0,1\}$ let $H_q(A,B)$ be the bipartite graph whose two vertex classes
are the difference lines and the sum lines of colour $q$. Its edges are the squares $(r,x)\in A\times B$
of colour $q$, the square $(r,x)$ joining the lines $r-x$ and $r+x$. For $z=(z_0,z_1)$ call $(A,B)$
\emph{$z$-feasible} if, for $q=0,1$, $H_q(A,B)$ has a vertex cover consisting of at most $z_q$
difference lines and at most $z_q$ sum lines.
\end{definition}

\begin{lemma}\label{lem:relax}
If $A$ is a set of empty rows and $B$ a set of empty columns of $D$, then $(A,B)$ is $(z_0,z_1)$-feasible.
\end{lemma}
\begin{proof}
A square $(r,x)\in A\times B$ is unoccupied, and no queen shares its row or its column. Hence it is
attacked along a diagonal, by a queen which then has the colour of $(r,x)$. So the difference and sum
lines of colour $q$ that carry a queen cover all edges of $H_q(A,B)$. There are at most $z_q$ of each
kind, since each queen of colour $q$ lies on one line of each kind and queens of colour $1-q$ lie on no
line of colour $q$.
\end{proof}

Feasibility is monotone: if $A'\subseteq A$ and $B'\subseteq B$, then $H_q(A',B')$ is a subgraph of
$H_q(A,B)$, so every sub-pair of a feasible pair is feasible. A search that builds $A$ and $B$ element by
element may therefore stop as soon as the partial pair is infeasible. Two certificates of
infeasibility are used.
\begin{enumerate}
\item[(M)] \emph{Matching bound.} A vertex cover of $H_q$ has at least $\nu(H_q)$ vertices, where $\nu$
is the matching number. So $\nu(H_q(A,B))>2z_q$ implies infeasibility. A matching of $H_q(A,B)$ is a set
of colour-$q$ squares of $A\times B$, no two on a common diagonal.
\item[(K)] \emph{Forced lines.} Distinct squares on one line lie on distinct lines of the other kind.
So a line containing more than $z_q$ squares of the graph belongs to every admissible cover, as
otherwise more than $z_q$ lines of the other kind would be needed. Removing forced lines, lowering the
budgets and iterating, and then applying (M) to the remaining graph against the remaining total budget,
gives a valid test.
\end{enumerate}

\begin{lemma}[symmetry]\label{lem:sym}
Let $n$ be even and $\bar X=\{n-1-i:i\in X\}$.
\begin{enumerate}
\item The transpose $(r,x)\mapsto(x,r)$, the half-turn $(r,x)\mapsto(n-1-r,n-1-x)$ and the anti-transpose
$(r,x)\mapsto(n-1-x,n-1-r)$ preserve colours. With the identity they form a group $K$ of order $4$. They
map a dominating set with empty rows $A$ and empty columns $B$ to a dominating set with the same colour
sizes and with empty rows and columns $(B,A)$, $(\bar A,\bar B)$ and $(\bar B,\bar A)$ respectively.
\item $z$-feasibility of $(A,B)$ is invariant under this action of $K$.
\item The reflection $(r,x)\mapsto(r,n-1-x)$ exchanges the two colours.
\end{enumerate}
\end{lemma}
\begin{proof}
(1) The colour of the image square is $x+r$, $2n-2-(r+x)$ and $2n-2-(r+x)$ respectively, all
$\equiv r+x$. The statement about empty lines is immediate. (2) The three maps send difference lines to
difference lines ($r-x\mapsto x-r$, $x-r$, $r-x$) and sum lines to sum lines ($r+x\mapsto r+x$,
$2n-2-(r+x)$, $2n-2-(r+x)$), preserving colours. Hence they induce isomorphisms $H_q(A,B)\cong
H_q(\sigma(A,B))$ which respect the two vertex classes, and admissible covers correspond. (3) The image
of $(r,x)$ has colour $r+n-1-x\equiv r+x+1$ because $n$ is even.
\end{proof}

\begin{proposition}\label{prop:complete}
Let $n$ be even and $g<n$. Let $\mathcal L$ be a set of pairs $(A,B,z_0)$, with $|A|=|B|=s$ and
$z_0\le\lfloor g/2\rfloor$, such that
\begin{enumerate}
\item[(a)] for every $z_0\le\lfloor g/2\rfloor$, every $(z_0,g-z_0)$-feasible pair of $s$-sets lies in
the $K$-orbit of a pair $(A,B)$ with $(A,B,z_0)\in\mathcal L$, and
\item[(b)] for every $(A,B,z_0)\in\mathcal L$, no set of $g$ queens contained in
$([0,n)\setminus A)\times([0,n)\setminus B)$ with at most $z_0$ queens of colour $0$ and at most
$g-z_0$ of colour $1$ dominates $Q_n$.
\end{enumerate}
Then $\gamma(Q_n)>g$.
\end{proposition}
\begin{proof}
Suppose $D$ is a dominating $g$-set. After the reflection of Lemma~\ref{lem:sym}(3) if needed,
$z_0(D)\le z_1(D)$, so $z_0(D)\le\lfloor g/2\rfloor$. Choose $s$-sets $A_0$, $B_0$ of empty rows and
columns. By Lemma~\ref{lem:relax}, $(A_0,B_0)$ is $(z_0(D),z_1(D))$-feasible, so by (a) there are
$\sigma\in K$ and $(A,B,z_0(D))\in\mathcal L$ with $\sigma(A_0,B_0)=(A,B)$. By Lemma~\ref{lem:sym}(1),
$\sigma D$ is a dominating $g$-set with the same colour sizes whose queens avoid the rows $A$ and the
columns $B$, contradicting~(b).
\end{proof}

\section{The computation}\label{sec:comp}

\paragraph{Stage 1 (pairs).} A depth-first search decides membership of rows and columns in $A$ and $B$
in the order row~$0$, column~$0$, row~$n{-}1$, column~$n{-}1$, row~$1$, column~$1$, \dots, from the
outside in. It maintains a maximum matching of $H_0\cup H_1$ incrementally, prunes with (M) and~(K),
and after every block of four decisions (rows and columns $t$ and $n-1-t$) discards prefixes that
cannot extend to the canonical member of their $K$-orbit. The canonical member is the one whose
sequence of block codes is lexicographically least, and the block code of an image depends only on the
same block of the original, so this test is exact. At the leaves it decides the constrained vertex
cover problem exactly, by branching. The output is a list of canonical feasible pairs. The search is
split into parts by the index of the search node at depth~$16$ modulo the number of parts. For $n=26$
there are $480$ parts per colour split; for $n=28$ there are $480$ parts for each $z_0\le4$ and $4\,800$
for each $z_0\ge5$.

\paragraph{Stage 2 (queens).} For each pair a branch-and-bound search finds all $g$-sets of queens
inside the allowed region, with the prescribed colour sizes, that dominate the whole board. It picks an
undominated square with the fewest available attackers and branches over those attackers, forbidding
earlier siblings so that the branches partition the search space. It prunes with greedy packings of
undominated squares whose sets of available attackers are pairwise disjoint, both overall and per colour.

\paragraph{The case $n=26$.} Here $g=s=13$ and $z_0\in\{0,\dots,6\}$. Table~\ref{tab:26} lists the work.
Stage~1 leaves canonical feasible pairs only for the split $(6,7)$, where there are $\clpairsSix$ of
them, and for the monochromatic split $(0,13)$. In the latter case there is exactly one pair: $A$ is the
set of odd rows and $B$ the set of even columns. This agrees with Weakley's structure
theory~\cite{Weakley2022}, which the computation does not use. Stage~2 finds no dominating set for any
of the $\clpairsTotal$ pairs. The SHA-256 hash of the survivor list begins with \texttt{\clsurvSha}.
In total stage~1 took $\clhStageOne$ and stage~2 $\clhStageTwo$ CPU hours (single cores of a shared
machine).

\begin{table}[ht]\centering\small
\begin{tabular}{lrrrrrrrr}\toprule
& \multicolumn{3}{c}{stage 1} & \multicolumn{3}{c}{stage 2} & \multicolumn{2}{c}{certificate}\\
\cmidrule(lr){2-4}\cmidrule(lr){5-7}\cmidrule(lr){8-9}
$(z_0,z_1)$ & nodes & pairs & CPU s & nodes & CPU s & sets & cubes & CPU s\\\midrule
\clrowZa
\clrowZb
\clrowZc
\clrowZd
\clrowZe
\clrowZf
\clrowZg
\bottomrule\end{tabular}
\caption{The computation for $n=26$, $g=13$. Stage~1 and stage~2 are the search programs of
Section~\ref{sec:comp}. The last two columns refer to the stage-1 certificate of
Section~\ref{sec:cert}: the number of solved cubes and the CPU time of CaDiCaL plus \texttt{drat-trim}.
Single-core CPU seconds on a shared machine.}\label{tab:26}
\end{table}

\paragraph{The case $n=28$.} Here $g=s=14$ and $z_0\in\{0,\dots,7\}$. Table~\ref{tab:28} lists the work
for the $\clpartsXXVIII$ parts. Stage~1 leaves canonical feasible pairs only for the splits $(7,7)$
($\clpairsSevenXXVIII$ pairs) and $(6,8)$ ($\clpairsSixXXVIII$ pairs), and for $(0,14)$ and $(1,13)$. In the
last two cases there is again exactly one pair, $A$ the odd rows and $B$ the even columns, whose grid
$A\times B$ consists of squares of a single colour. This again agrees with Weakley's results: by his
Theorem~17 and Proposition~11~\cite{Weakley2022} a $14$-dominating set of $Q_{28}$ would have colour
classes $\{7,7\}$ or $\{6,8\}$. Altogether there are $\clpairsTotalXXVIII$ pairs; the SHA-256 hash of the
survivor list begins with \texttt{\clsurvShaXXVIII}. Stage~1 took $\clhStageOneXXVIII$ CPU hours. Of the
parts, $\clcloudPartsXXVIII$ (all but $16$ of those with $z_0=7$, and $2673$ of those with $z_0=6$) were run on
rented single-core cloud containers (Modal), compiled there from the same source inside a pinned
Debian image. The cost was about US\$4. Every returned part was checked locally against its log and its
SHA-256 hash before it was accepted, and the remaining parts, including $16$ that exceeded the time limit
of the cloud calls, ran locally. The cloud cores were about $1.8$ times slower than the local ones, which
inflates the CPU times of Table~\ref{tab:28}.

For $n=28$ condition~(b) of Proposition~\ref{prop:complete} rests on the SAT check of
Section~\ref{sec:cert} alone. The branch-and-bound search of stage~2 was run only on the
$\clplacePairsXXVIII$ pairs with $z_0\in\{1,6\}$, in $\clplaceCpuXXVIII$\,s; it found no dominating set.
Its running time has a heavy tail (up to $10^3$\,s for a single pair with $z_0=6$), so we did not run it
on the $\clpairsSevenXXVIII$ pairs with $z_0=7$. On the pair with $z_0=0$ we stopped it after $3.6$ hours.

\begin{table}[ht]\centering\small
\begin{tabular}{lrrrrrr}\toprule
& \multicolumn{3}{c}{stage 1} & \multicolumn{2}{c}{certificate} & \multicolumn{1}{c}{SAT stage 2}\\
\cmidrule(lr){2-4}\cmidrule(lr){5-6}\cmidrule(lr){7-7}
$(z_0,z_1)$ & nodes & pairs & CPU s & cubes & CPU s & CPU s\\\midrule
\clrowYa
\clrowYb
\clrowYc
\clrowYd
\clrowYe
\clrowYf
\clrowYg
\clrowYh
\bottomrule\end{tabular}
\caption{The computation for $n=28$, $g=14$. Stage~1 is the search program of Section~\ref{sec:comp};
most of the stage-1 time for $(7,7)$ and part of it for $(6,8)$ was spent on slower cloud cores. The
certificate columns refer to the stage-1 certificate of Section~\ref{sec:cert}, and the last column to
the SAT check of condition~(b), each CaDiCaL plus \texttt{drat-trim}. Single-core CPU seconds.}\label{tab:28}
\end{table}

\section{Certificates}\label{sec:cert}

The two searches of Section~\ref{sec:comp} are fast but intricate programs. Both conditions of
Proposition~\ref{prop:complete} were therefore also established by SAT solving with proofs: CaDiCaL
\clcadicalVersion~\cite{CaDiCaL} writes a DRAT proof of unsatisfiability for each formula and
\texttt{drat-trim}~\cite{DRATtrim} checks it. The formulas are produced by two short encoders that do not
share code with the search programs. Cardinality constraints use the sequential counter~\cite{Sinz}.

\paragraph{Condition (b).} For each listed pair, the formula has one variable per allowed square, one
clause per square of the board stating that some allowed square attacking it (or equal to it) carries a
queen, and constraints that at most $z_q$ queens have colour~$q$. For $n=26$ all $\clpairsTotal$
formulas are unsatisfiable, and all $\clpairsTotal$ proofs were verified. CaDiCaL took $\clsatSolve$\,s
in total and \texttt{drat-trim} $\clsatCheck$\,s. The monochromatic pair is by far the hardest instance
($\clsatMonoSolve$\,s and $\clsatMonoCheck$\,s); each other instance takes less than $\clsatOtherMax$\,s for
each of the two steps. For $n=28$ all $\clpairsTotalXXVIII$ formulas are unsatisfiable and all proofs were
verified, in $\clsatSolveXXVIII$\,s (CaDiCaL) and $\clsatCheckXXVIII$\,s (\texttt{drat-trim}). The hardest
instance is again the pair whose grid has a single colour, here with $z_0=0$ ($\clsatMonoSolveXXVIII$\,s and
$\clsatMonoCheckXXVIII$\,s); every other instance takes less than $\clsatOtherMaxXXVIII$\,s for each step.

\paragraph{Condition (a).} Fix $z_0$. The formula $F_{z_0}$ has variables $a_r$ (row $r$ is in $A$), $b_x$
(column $x$ is in $B$), $y_d$ (difference line $d$ in the cover) and $w_t$ (sum line $t$ in the cover), with the
constraints $|A|=|B|=s$, at most $z_q$ chosen difference lines and at most $z_q$ chosen sum lines of
colour~$q$, and, for every square, the clause $\neg a_r\vee\neg b_x\vee y_{r-x}\vee w_{r+x}$. The
projections of its models to the variables $a,b$ are exactly the $z$-feasible pairs. For each listed
pair $(A,B,z_0)$ and each of its four images $(P,Q)$ under $K$ we add the clause
$\bigvee_{r\in P}\neg a_r\vee\bigvee_{x\in Q}\neg b_x$, which excludes $(P,Q)$. Condition~(a) for this
$z_0$ says exactly that the resulting formula $F'_{z_0}$ is unsatisfiable.

We refute $F'_{z_0}$ by cube and conquer~\cite{CubeConquer}, with a symmetry reduction. Fix the $4T$
variables $a_t,a_{n-1-t},b_t,b_{n-1-t}$ ($t<T$) in all $16^T$ ways; these cubes cover all assignments.
The group $K$ acts on the cubes, and $F'_{z_0}\wedge c$ is satisfiable if and only if
$F'_{z_0}\wedge\sigma c$ is: the projections of the models of $F'_{z_0}$ form a $K$-invariant set of
pairs, by Lemma~\ref{lem:sym}(2) and because the excluded set is a union of orbits. So it suffices to
refute $F'_{z_0}\wedge c$ for one cube $c$ in each orbit. We used $T=2$ ($\clcubesSmall$ orbits) for
$z_0\le5$ and $T=3$ ($\clcubesSix$ orbits) for $z_0=6$ and, for $n=28$, $z_0=7$. A short script checks
that the orbits of the chosen cubes cover all $16^T$ cubes. Every cube formula was refuted and every
proof verified. For $n=26$ this took $\clhCert$ CPU hours and the proofs were about $\clcertProofGB$\,GB in
all. For $n=28$ it took $\clhCertXXVIII$ CPU hours; for $z_0=7$ the formula has $\clclausesSevenXXVIII$
clauses, most of them exclusion clauses, the largest proof was $\clcertMaxProofGBXXVIII$\,GB and all proofs
together about $\clcertProofGBXXVIII$\,GB. The proofs were deleted after checking. The encoder rebuilds the
formulas deterministically, their SHA-256 hashes are recorded in the repository, and the verification
script regenerates and checks a random sample of the proofs. For every $z_0\ge3$ and both $n$ the
certificate took less CPU time than the stage-1 search itself (Tables~\ref{tab:26} and~\ref{tab:28}).

\paragraph{What the proof rests on.} Theorem~\ref{thm:main} follows from Proposition~\ref{prop:complete},
applied for $n=26$ and for $n=28$ with $\mathcal L$ the respective survivor list, and from the upper bounds
of~\cite{OW,OEIS} for $n=26$ and of~\cite{LMR} and Section~\ref{sec:valid} for $n=28$. Its correctness depends on
Lemmas~\ref{lem:relax} and~\ref{lem:sym} and Proposition~\ref{prop:complete}, which are proved above;
on the two CNF encoders (\texttt{relax\_cnf} in \texttt{code/refcheck.py}, about 30 lines using the
cardinality encodings of PySAT~\cite{PySAT}, and \texttt{build} in \texttt{code/sat\_place.py}, about 40
lines); on the short enumeration of cube orbit representatives and its check that their
orbits cover all $16^T$ cubes; and on \texttt{drat-trim}. It does not depend on the correctness of the
two search programs, which serve only to produce the lists $\mathcal L$; in particular it does not depend
on the cloud computations. As a check that the encoders are not too strong, every listed pair, for both
$n$, satisfies $F_{z_0}$ without the exclusion clauses, the two known $14$-queen dominating sets of
$Q_{26}$ satisfy the stage-2 encoding for $g=14$, and the $15$-queen dominating set of $Q_{28}$ below
satisfies it for $g=15$.

\section{Validation and independent checks}\label{sec:valid}

\paragraph{Known values.} Table~\ref{tab:valid} shows the pipeline for $g=n/2$ and even $n\le18$. The
numbers of $n/2$-dominating sets up to the dihedral symmetries agree with the
literature~\cite{Bird,RB}, and for $n=8,14,16$ no $n/2$-set exists. The recorded runs for $n=20,22,24$
($\clpairsrecTwenty$, $\clpairsrecTwentyTwo$ and $\clpairsrecTwentyFour$ feasible pairs) found no
$n/2$-set, in agreement with~\cite{Bird}; the verification script reruns $n=20,22$ in its full mode.

\begin{table}[ht]\centering
\begin{tabular}{lrrrrrrrrrrrrr}\toprule
$n$ & 4&6&8&10&12&14&16&18&20&22&24&26&28\\\midrule
number of $n/2$-sets up to symmetry & \clvalFour&\clvalSix&\clvalEight&\clvalTen&\clvalTwelve&\clvalFourteen&\clvalSixteen&\clvalEighteen&0&0&0&\textbf{0}&\textbf{0}\\
\bottomrule\end{tabular}
\caption{Numbers of dominating sets of size $n/2$ of $Q_n$ up to the dihedral group, as found by the
pipeline. The values for $n\le 24$ agree with~\cite{Bird,RB}; those for $n=26,28$ are Theorem~\ref{thm:main}.}\label{tab:valid}
\end{table}

All $n/2$-sets in Table~\ref{tab:valid} with $n>4$ are monochromatic. To test the pipeline on
bichromatic sets, we also ran it with $g=\gamma(Q_n)=n/2+1$ for $n=8,14,16$, where minimum dominating
sets are plentiful and mostly bichromatic. Stage~2 examined $\clpairsgpEight$, $\clpairsgpFourteen$ and
$\clpairsgpSixteen$ feasible pairs and found $\clgpEight$, $\clgpFourteen$ and $\clgpSixteen$ classes. These
agree with~\cite{RB}, and with~\cite{Bird} for $n=8,14$; for $n=16$ Bird lists $43$ classes, which
Rostami and Bright corrected to $371$.

\paragraph{Second implementations.} Stage~1 was reimplemented separately. The second program uses
another decision order, breadth-first augmenting paths, only bound~(M), another exact cover algorithm
and another canonical form. For every $n\le\clrefCompareMaxN$ and for $n=26$ the two programs produce
the same set of orbits of feasible pairs; for $n=26$ the second program needed $\clvNodes$ nodes and
$\clvCpu$\,s. For $n\le\clbruteMaxN$ both also agree with a test of every pair $(A,B)$, without pruning
or symmetry. Stage~2 has the SAT check described above as its second implementation.

For $n=28$ the second implementations were run only in part, since the certificates make the lower
bound independent of all search programs. The second stage-1 program was run on every part with
$z_0\le5$, where it finds the same orbits as stage~1 (the two single-colour pairs and nothing else), and
on $\clvPartsSixXXVIII$ of the $4\,800$ parts with $z_0=6$, where its $\clvPairsSixXXVIII$ pairs are all among
those found by stage~1; it needed $\clvNodesXXVIII$ nodes and $\clvCpuXXVIII$\,s. It was not run for
$z_0=7$. The branch-and-bound search of stage~2 was run on the pairs with $z_0\in\{1,6\}$, as described in
Section~\ref{sec:comp}.

\paragraph{Independent referee checks.} An independent AI referee, not involved in writing the search
programs, reviewed the proofs of Section~\ref{sec:relax} and the two stage-1 programs. It also wrote
the program \texttt{code/refcheck.py}, which uses only SAT solving and shares no code with the search
programs. This program enumerates all feasible pairs by incremental SAT solving with orbit-blocking clauses and finds the same
orbit sets as stage~1 for every even $n\le\clrefCompareMaxN$ and every $z_0$. It runs a complete
independent pipeline (SAT enumeration of pairs, then SAT enumeration of all dominating sets inside each
region), which reproduces the class counts above for $n\le\clrefPipeMaxN$ and for $(n,g)=(8,5)$. It also
enumerates all $g$-dominating sets directly, without the relaxation, in the cases with
$n\le\clrefDirectMaxN$, again with the same counts. The referee also proposed the stage-1 certificate of
Section~\ref{sec:cert}; the same certificate is produced in the verification run for $n=16$ and $n=18$.

\paragraph{Upper bounds.} The verification script checks that Kamenetsky's $14$-queen set and the
independent $14$-queen set of \"Osterg\aa rd and Weakley dominate $Q_{26}$, and that their $14$-queen
set of Corollary~\ref{cor:oeis} dominates $Q_{27}$. For $n=28$ we found a $15$-queen dominating set by
simulated annealing; it is in the repository, and the script checks it. (The value $i(Q_{28})=15$
of~\cite{LMR} already gives $\gamma(Q_{28})\le15$.) The script also checks the sets used for $n=29,30,31$
in the proof of Corollary~\ref{cor:oeis}.

\section{Remarks}

The relaxation of Lemma~\ref{lem:relax} keeps only the line structure of the squares that must be
covered diagonally, yet it is strong enough that stage~1 leaves only $\clpairsTotal$ candidate pairs for
$n=26$ and $\clpairsTotalXXVIII$ for $n=28$. The cost grows quickly: from $n=26$ to $n=28$ the stage-1
search grew from $\clhStageOne$ to $\clhStageOneXXVIII$ CPU hours (part of the latter on slower cores), the
number of pairs by a factor of about $11$, and the stage-1 certificate from $\clhCert$ to $\clhCertXXVIII$ CPU
hours. The smallest open case is now $n=32$. There $4\mid n$, so again only bichromatic $16$-sets are
possible~\cite{Weakley2022}. Extrapolating the growth from $n=26$ to $n=28$ over two more steps suggests
some $10^3$ to $10^4$ CPU hours for the present method, so deciding $\gamma(Q_{32})\in\{16,17\}$ probably
calls for a stronger relaxation, or for using structural results such as Proposition~11
of~\cite{Weakley2022} as proved lemmas.

\section*{Reproducibility}
The repository\footnote{\repo} contains the C sources of both stages and of the second stage-1 program,
the CNF encoders, the referee's checks, the driver used for the cloud runs, and for $n=26$ and $n=28$ the
survivor lists, per-part SHA-256 manifests, the logs of the searches, the per-pair SAT records and the
per-cube records of the stage-1 certificates. The script \texttt{scripts/verify.py} rebuilds everything
and, in about twenty minutes on three cores, reruns the validation for small $n$, the referee's checks
and the stage-1 certificate for $n=16,18$; checks the records and their hashes for $n=26$ and $n=28$,
including the base formulas of both certificates, rebuilt from the survivor lists; checks the upper
bounds; and reruns seeded random samples of the two computations against the records ($\clsampleParts$
and $\clsamplePartsXXVIII$ stage-1 parts, SAT instances and certificate cubes). Every number quoted in this
paper is generated from a file of claims that the script compares with the computed values. With
\texttt{--full} it reruns the whole $n=26$ computation, and with \texttt{--heavy28} a larger part of the
$n=28$ computation (about six CPU hours).

\section*{Declaration of AI-generated research and authorship}
This manuscript is the product of an AI-driven research process. During a broader investigation of
open problems, initiated by the human author, an AI system identified this problem as a candidate,
devised the method, wrote and ran the code for the computations, and wrote this manuscript. AI systems
produced the results and their presentation, including the mathematical exposition, the computations
and the organisation and writing of the text; the human author directed the AI systems but wrote no
part of the manuscript. The system was Claude Opus~5.5 (Anthropic), used through Claude Code.

The computational claims are reproduced by \texttt{scripts/verify.py}. At the time of writing, the
mathematics and the code have been checked only by AI systems and by computer verification. These
checks include an adversarial review by an independent AI referee, who re-derived the proofs of
Section~\ref{sec:relax}, read the search programs, wrote the independent SAT-based checks of
Section~\ref{sec:valid} and the stage-1 certificate of Section~\ref{sec:cert}. No independent human
expert has reviewed the work. The human author has not read this manuscript or checked its arguments.
The author takes responsibility for the accuracy of this disclosure.

\raggedright
\begin{thebibliography}{99}
\bibitem{CaDiCaL} A.~Biere, T.~Faller, K.~Fazekas, M.~Fleury, N.~Froleyks, F.~Pollitt, \emph{CaDiCaL~2.0},
in: Computer Aided Verification (CAV 2024), Lecture Notes in Comput.\ Sci., Springer, 2024, 133--152.
\href{https://doi.org/10.1007/978-3-031-65627-9_7}{doi:10.1007/978-3-031-65627-9\_7}.
\bibitem{Bird} W.~H.~Bird, \emph{Computational methods for domination problems}, PhD thesis, University of
Victoria, 2017. \url{https://dspace.library.uvic.ca/handle/1828/8634}.
\bibitem{Cockayne} E.~J.~Cockayne, \emph{Chessboard domination problems}, Discrete Math.\ \textbf{86}
(1990), 13--20. \href{https://doi.org/10.1016/0012-365X(90)90344-H}{doi:10.1016/0012-365X(90)90344-H}.
\bibitem{FW} D.~Finozhenok, W.~D.~Weakley, \emph{An improved lower bound for domination numbers of the
queen's graph}, Australas.\ J.\ Combin.\ \textbf{37} (2007), 295--300.
\bibitem{CubeConquer} M.~J.~H.~Heule, O.~Kullmann, S.~Wieringa, A.~Biere, \emph{Cube and conquer:
guiding CDCL SAT solvers by lookaheads}, in: Hardware and Software: Verification and Testing (HVC 2011),
Lecture Notes in Comput.\ Sci., Springer, 2012, 50--65.
\href{https://doi.org/10.1007/978-3-642-34188-5_8}{doi:10.1007/978-3-642-34188-5\_8}.
\bibitem{PySAT} A.~Ignatiev, A.~Morgado, J.~Marques-Silva, \emph{PySAT: a Python toolkit for prototyping
with SAT oracles}, in: Theory and Applications of Satisfiability Testing (SAT 2018), Lecture Notes in
Comput.\ Sci., Springer, 2018, 428--437.
\href{https://doi.org/10.1007/978-3-319-94144-8_26}{doi:10.1007/978-3-319-94144-8\_26}.
\bibitem{Kearse} M.~D.~Kearse, P.~B.~Gibbons, \emph{Computational methods and new results for chessboard
problems}, Australas.\ J.\ Combin.\ \textbf{23} (2001), 253--284.
\bibitem{LMR} A.~Langlois-R\'emillard, M.~M\"ussig, \'E.~Rold\'an, \emph{Complexity of chess domination
problems}, Res.\ Math.\ Sci.\ \textbf{12} (2025), no.~1, paper no.~14.
\href{https://doi.org/10.1007/s40687-024-00492-5}{doi:10.1007/s40687-024-00492-5}.
\bibitem{OEIS} OEIS Foundation Inc., \emph{The On-Line Encyclopedia of Integer Sequences}, sequence
A075458, with D.~Kamenetsky's file of best known placements for $n\le26$ (2019),
\url{https://oeis.org/A075458}.
\bibitem{OW} P.~R.~J.~\"Osterg\aa rd, W.~D.~Weakley, \emph{Values of domination numbers of the queen's graph},
Electron.\ J.\ Combin.\ \textbf{8} (2001), \#R29.
\href{https://doi.org/10.37236/1573}{doi:10.37236/1573}.
\bibitem{RV} V.~Raghavan, S.~M.~Venkatesan, \emph{On bounds for a board covering problem}, Inform.\ Process.\
Lett.\ \textbf{25} (1987), 281--284.
\href{https://doi.org/10.1016/0020-0190(87)90201-8}{doi:10.1016/0020-0190(87)90201-8}.
\bibitem{RB} T.~Rostami, C.~Bright, \emph{Queen domination by SAT solving}, preprint (2025),
arXiv:2508.11945. \href{https://doi.org/10.48550/arXiv.2508.11945}{doi:10.48550/arXiv.2508.11945}.
\bibitem{Sinz} C.~Sinz, \emph{Towards an optimal CNF encoding of Boolean cardinality constraints}, in:
Principles and Practice of Constraint Programming (CP 2005), Lecture Notes in Comput.\ Sci., Springer, 2005, 827--831.
\href{https://doi.org/10.1007/11564751_73}{doi:10.1007/11564751\_73}.
\bibitem{Watkins} J.~J.~Watkins, \emph{Across the Board: The Mathematics of Chessboard Problems}, Princeton
University Press, 2004. \href{https://doi.org/10.1515/9781400840922}{doi:10.1515/9781400840922}.
\bibitem{Weakley2002} W.~D.~Weakley, \emph{A lower bound for domination numbers of the queen's graph},
J.\ Combin.\ Math.\ Combin.\ Comput.\ \textbf{43} (2002), 231--254.
\bibitem{Weakley2022} W.~D.~Weakley, \emph{Queen domination of even square boards}, Electron.\ J.\ Combin.\
\textbf{29}(2) (2022), \#P2.50.
\href{https://doi.org/10.37236/10617}{doi:10.37236/10617}.
\bibitem{DRATtrim} N.~Wetzler, M.~J.~H.~Heule, W.~A.~Hunt~Jr., \emph{DRAT-trim: efficient checking and
trimming using expressive clausal proofs}, in: Theory and Applications of Satisfiability Testing (SAT 2014),
Lecture Notes in Comput.\ Sci., Springer, 2014, 422--429.
\href{https://doi.org/10.1007/978-3-319-09284-3_31}{doi:10.1007/978-3-319-09284-3\_31}.
\end{thebibliography}

\end{document}
